\SourceLeaf{052}%
veris discrepare censebuntur, ita ut ipsorum differentias a valoribus observatis eo maiore iure tamquam observationum errores veros adoptare liceat, quo maior fuerit harum multitudo. Hanc praxin sequuti sunt omnes calculatores, qui observationum praecisionem in casibus concretis a posteriori aestimare susceperunt: sed manifesto illa theoretice erronea est, et quamquam in casibus multis ad usus practicos sufficere possit, tamen in aliis enormiter peccare potest. Summopere itaque hoc argumentum dignum est, quod accuratius enodetur.

Retinebimus in hac disquisitione designationes inde ab art.\ 19 adhibitas. Praxis ea, de qua diximus, quantitates \(A,B,C\) etc. tamquam valores veros ipsarum \(x,y,z\) considerat, et proin ipsas \(\lambda,\lambda',\lambda''\) etc. tamquam valores veros functionum \(v,v',v''\) etc. Si omnes observationes aequali praecisione gaudent, ipsarumque pondus \(p=p'=p''\) etc. pro unitate acceptum est, eaedem quantitates, signis mutatis, in illa suppositione observationum errores exhibent, unde praecepta art.\ 15 praebent observationum errorem medium \(m\)
\[
=\BareRoot{\frac{\lambda\lambda+\lambda'\lambda'+\lambda''\lambda''+\text{etc.}}{\pi}}=\BareRoot{\frac{M}{\pi}}
\]
Si observationum praecisio non est eadem, quantitates \(-\lambda,-\lambda',-\lambda''\) etc. exhibent observationum errores per radices quadratas e ponderibus multiplicatos, praeceptaque art.\ 16 ad eandem formulam \(\BareRoot{\frac{M}{\pi}}\) perducerent, iam errorem medium talium observationum, quibus pondus \(=1\) tribuitur, denotantem. Sed manifesto calculus exactus requireret, ut loco quantitatum \(\lambda,\lambda',\lambda''\) etc. valores functionum \(v,v',v''\) etc. e valoribus veris ipsarum \(x,y,z\) etc. prodeuntes adhiberentur, i.\ e. loco ipsius \(M\), valor functionis \(\Omega\) valoribus veris ipsarum \(x,y,z\) etc. respondens. Qui quamquam assignari nequeat, tamen certi sumus, eum esse maiorem quam \(M\) (quippe qui est minimus possibilis), excipiendo casum infinite parum probabilem, ubi incognitarum valores maxime plausibiles exacte cum veris quadrant. In genere itaque affirmare possumus, praxin vulgarem errorem medium iusto minorem producere, sive observationibus praecisionem nimis magnam tribuere. Videamus iam, quid doceat theoria rigorosa.

\Article{38}
Ante omnia investigare oportet, quonam modo \(M\) ab observationum erroribus veris pendeat. Denotemus hos, ut in art.\ 28, per \(e,e',e''\) etc., statuamusque ad maiorem simplicitatem
\SourceLeaf{053}
\[
\begin{gathered}
e\BareRoot{p}=\varepsilon,\quad e'\BareRoot{p'}=\varepsilon',\quad e''\BareRoot{p''}=\varepsilon''\quad\text{etc., nec non}\\
m\BareRoot{p}=m'\BareRoot{p'}=m''\BareRoot{p''}\quad\text{etc.}=\mu
\end{gathered}
\]
Porro sint valores veri ipsarum \(x,y,z\) etc., resp. \(A-x^0,\allowbreak B-y^0,\allowbreak C-z^0\) etc., quibus respondeant valores ipsarum \(\xi,\eta,\zeta\) etc. hi \(-\xi^0,-\eta^0,-\zeta^0\) etc. Manifesto iisdem respondebunt valores ipsarum \(v,v',v''\) etc. hi \(-\varepsilon,-\varepsilon',-\varepsilon''\) etc., ita ut habeatur
\[
\begin{aligned}
\xi^0&=a\varepsilon+a'\varepsilon'+a''\varepsilon''+\text{etc.}\\
\eta^0&=b\varepsilon+b'\varepsilon'+b''\varepsilon''+\text{etc.}\\
\zeta^0&=c\varepsilon+c'\varepsilon'+c''\varepsilon''+\text{etc.}
\end{aligned}
\]
etc. nec non
\[
\begin{aligned}
x^0&=\alpha\varepsilon+\alpha'\varepsilon'+\alpha''\varepsilon''+\text{etc.}\\
y^0&=\beta\varepsilon+\beta'\varepsilon'+\beta''\varepsilon''+\text{etc.}\\
z^0&=\gamma\varepsilon+\gamma'\varepsilon'+\gamma''\varepsilon''+\text{etc.}
\end{aligned}
\]
Denique statuamus
\[
\Omega^0=\varepsilon\varepsilon+\varepsilon'\varepsilon'+\varepsilon''\varepsilon''+\text{etc.}
\]
ita ut sit \(\Omega^0\) aequalis valori functionis \(\Omega\), valoribus veris ipsarum \(x,y,z\) etc. respondenti. Hinc quum habeatur indefinite
\[
\Omega=M+(x-A)\xi+(y-B)\eta+(z-C)\zeta+\text{etc.}
\]
erit etiam
\[
M=\Omega^0-x^0\xi^0-y^0\eta^0-z^0\zeta^0-\text{etc.}
\]
Hinc manifestum est, \(M\), evolutione facta esse functionem homogeneam secundi ordinis errorum \(e,e',e''\) etc., quae, pro diversis errorum valoribus maior minorve evadere poterit. Sed quatenus errorum magnitudo nobis incognita manet, functionem hanc indefinite considerare, imprimisque secundum principia calculi probabilitatis eius valorem medium assignare conveniet. Quem inveniemus, si loco quadratorum \(ee,e'e',e''e''\) etc. resp. scribimus \(mm,m'm',m''m''\) etc., producta vero \(ee',ee'',e'e''\) etc omnino omittimus, vel quod idem est, si loco cuiusvis quadrati \(\varepsilon\varepsilon,\varepsilon'\varepsilon',\varepsilon''\varepsilon''\) etc. scribimus \(\mu\mu\), productis \(\varepsilon\varepsilon',\varepsilon\varepsilon'',\varepsilon'\varepsilon''\) etc. prorsus neglectis. Hoc modo e termino \(\Omega^0\) manifesto provenit \(\pi\mu\mu\); terminus \(-x^0\xi^0\) producet
\[
-(a\alpha+a'\alpha'+a''\alpha''+\text{etc.})\mu\mu=-\mu\mu
\]
\SourceLeaf{054}%
et similiter singulae partes reliquae praebebunt \(-\mu\mu\), ita ut valor medius totalis fiat \(=(\pi-\rho)\mu\mu\), denotante \(\pi\) multitudinem observationum, \(\rho\) multitudinem incognitarum. Valor verus quidem ipsius \(M\), prout fors errores obtulit, maior minorve medio fieri potest, sed discrepantia eo minoris momenti erit, quo maior fuerit observationum multitudo, ita ut pro valore approximato ipsius \(\mu\) accipere liceat
\[
\BareRoot{\frac{M}{\pi-\rho}}
\]
Valor itaque ipsius \(\mu\), ex praxi erronea, de qua in art. praec. loquuti sumus, prodiens, augeri debet in ratione quantitatis \(\BareRoot{(\pi-\rho)}\) ad \(\BareRoot{\pi}\).

\Article{39}
Quo clarius eluceat, quanto iure valorem fortuitum ipsius \(M\) medio aequiparare liceat, adhuc investigare oportet errorem medium metuendum, dum statuimus \(\frac{M}{\pi-\rho}=\mu\mu\). Iste error medius aequalis est radici quadratae e valore medio quantitatis
\[
\left(\frac{\Omega^0-x^0\xi^0-y^0\eta^0-z^0\zeta^0-\text{etc.}-(\pi-\rho)\mu\mu}{\pi-\rho}\right)^2
\]
quam ita exhibebimus
\[
\begin{aligned}
&\left(\frac{\Omega^0-x^0\xi^0-y^0\eta^0-z^0\zeta^0-\text{etc.}}{\pi-\rho}\right)^2\\[4pt]
&\quad-\frac{2\mu\mu}{\pi-\rho}(\Omega^0-x^0\xi^0-y^0\eta^0-z^0\zeta^0-\text{etc.}-(\pi-\rho)\mu\mu)-\mu^4
\end{aligned}
\]
et quum manifesto valor medius termini secundi fiat \(=0\), res in eo vertitur, ut indagemus valorem medium functionis
\[
\Psi=(\Omega^0-x^0\xi^0-y^0\eta^0-z^0\zeta^0-\text{etc.})^2
\]
quo invento et per \(N\) designato, error medius quaesitus erit
\[
=\BareRoot{\left(\frac{N}{(\pi-\rho)^2}-\mu^4\right)}
\]
Expressio \(\Psi\) evoluta manifesto est functio homogenea sive errorum \(e,e',e''\) etc., sive quantitatum \(\varepsilon,\varepsilon',\varepsilon''\) etc., eiusque valor medius invenietur, si

\(1^\mathrm{o}\) pro biquadratis \(e^4,e'^4,e''^4\) etc. substituuntur eorum valores medii

\(2^\mathrm{o}\) pro singulis productis e binis quadratis ut \(eee'e',\allowbreak eee''e'',\allowbreak e'e'e''e''\) etc. producta ex ipsorum valoribus mediis, puta \(mmm'm',\allowbreak mmm''m'',\allowbreak m'm'm''m''\) etc.
